% ===== 原文 pp.12--24 =====
% 块边界说明：p.12 顶部两行是 §1 末段的跨页尾巴（该句起于 p.11），按 BRIEF「跨页句子由开始者写完」
% 已由 Lead 写入 c01.tex；本块自 p.12 的 Part I 标题起。

% [原文 p.12]
\part{第一篇\quad Class S 理论}

\section{ADE 型的 6d $(2,0)$ SCFT}

超共形代数只存在于至多 6 维的时空中；如今已有大量证据表明存在 6d $\cN=(2,0)$（最大超对称）的 SCFT $\cX(\fg)$，它由单链（simply-laced）李代数 $\fg$ 标记。\footnote{提醒一下：单链李代数是 $\mathfrak{a}_{N-1}=\fsu(N)$、$\mathfrak{so}(2N)=\mathfrak{d}_N$，以及三个例外代数 $\mathfrak{e}_6,\mathfrak{e}_7,\mathfrak{e}_8$（每种情形的下标都是秩）。这一 ADE 分类在理论物理中还有若干优美的化身，但本综述不打算展开。}\textsuperscript{,}\footnote{虽然对 $\cX(\fg)$ 的不同构造给出的是同一个条件，即 $\fg$ 为单链（其中也包括一些场论论证 [152]），但 $\cX(\fg)$ 是否穷尽了所有 6d $\cN=(2,0)$ SCFT 尚未被证明。4d $\cN=4$ SCFT 的情形相同：可能存在 $\cN=4$ SYM 理论之外的这类理论。} 没有人知道如何真正用 QFT 语言------例如通过拉格朗日表述------来定义 $\cX(\fg)$。它只能作为某些弦论或 M 理论膜构型的退耦极限而得到。

% [Tutor Note: p.12 的要点框在原文中是右缘绕排浮框（正文被压成窄栏、框贴右边），
%  因此英文原文说 "properties detailed to the right"。译文改为常规单栏排版，框紧跟正文之后，
%  故正文相应改说「详见下方要点框」。核对方法：p.12.png 左侧五行正文的行宽明显短于正文栏宽，
%  且 "Despite its stringy construction" 一段的右边界正好贴着灰框左边。]

尽管构造是弦论的，该理论仍被期望是一个名副其实的局域量子场论，例如具有局域守恒的能动张量。\footnote{准确地说，它是一个相对量子场论（relative quantum field theory）[153]。} 这些构造蕴含三条重要性质，详见下方要点框。

\begin{keybox}{$\cX(\fg)$ 的性质}
\begin{itemize}[leftmargin=1.4em,itemsep=1pt,topsep=1pt]
\item $\cX(\fg)$ 的真空，其红外（IR）描述是自对偶二形式规范场的阿贝尔 6d $(2,0)$ 理论，这些场取值于 $\fg$ 的 Cartan 代数、并模掉 Weyl 群。
\item $\cX(\fg)$ 是 5d $\cN=2$ SYM 的紫外完备，意思是：规范代数为 $\fg$、规范耦合为 $g_{5d}$ 的 SYM，给出 $\cX(\fg)$ 紧化在半径为 $g_{5d}^2$ 的圆上的红外描述。
\item $\cX(\fg)$ 允许由 $\fg$ 中幂零轨道（nilpotent orbit）标记的余维 2 半 BPS 缺陷，以及由 $\fg$ 的表示标记的余维 4 半 BPS 缺陷。
\end{itemize}
\end{keybox}

前两条性质彼此相容，因为 5d $\cN=2$ SYM 在其 Coulomb 分支上、以及圆上的阿贝尔 6d 理论，都由 $\fg$ 的 Cartan 中取值的 5d 阿贝尔矢量多重态描述。最后一条性质也同样相容：当沿 Coulomb 分支移动，或把理论放到圆上时，这些缺陷都有相当显式的描述。具有这些性质的 $\cX(\fg)$ 的存在性，已由许多涉及更易理解理论的相容性检验所证实。其中一大类相容性检验，就是把这些理论放到 4 维流形与 2 维流形的乘积 $M_4\times C_2$ 上所得到的 AGT 对应。

% [原文 p.13]
\begin{table}[htbp]
\centering
\small
\begin{tabular}{lllll}
\toprule
 & 超代数 & 共形代数 & R 对称代数 & $Q\&S$ \\
\midrule
3d $\cN\le 8$ & $\mathfrak{osp}(\cN|4)$ & $\fsp(4,\R)$ & $\fso(\cN)$ & $(\mathbf{4},\cN)$ \\
4d $\cN\le 3$ & $\fsu(2,2|\cN)$ & $\fsu(2,2)$ & $\fsu(\cN)\oplus\fu(1)$ & $(\mathbf{4},\overline{\cN})\oplus(\overline{\mathbf{4}},\cN)$ \\
4d $\cN=4$ & $\mathfrak{psu}(2,2|4)$ & $\fsu(2,2)$ & $\fsu(4)$ & $(\mathbf{4},\overline{\mathbf{4}})\oplus(\overline{\mathbf{4}},\mathbf{4})$ \\
5d $\cN=1$ & $\mathfrak{f}^2(4)$ & $\fso(2,5)$ & $\fsu(2)$ & $(\mathbf{8},\mathbf{2})$ \\
6d $\cN\le 2$ & $\mathfrak{osp}(8^*|2\cN)$ & $\fso^*(8)$ & $\mathfrak{usp}(2\cN)$ & $(\mathbf{8},\mathbf{2}\cN)$ \\
\bottomrule
\end{tabular}
\tabtitle{2}{洛伦兹号差下超共形代数的 Nahm 分类。这里列出各维数下的超共形代数、两个玻色因子（共形代数与 R 对称代数），以及庞加莱超荷 $Q$ 与共形超荷 $S$ 在其中变换的表示（指这些玻色因子的表示）。}
\end{table}

% [Tutor Note: p.13 表 2 的 5d 一行，原文把该超代数排作 $\mathfrak{f}^2(4)$（上标 2 在 fraktur f 之后）。
%  按通常记号它应写作 $F(4)$／$\mathfrak{f}(4)$：例外李超代数，偶部为 $\fso(2,5)\oplus\fsu(2)$（正是表中
%  "Conformal" 与 "R-symmetry" 两列），共 8 个超荷。核对方法：数 "Q&S" 列维数 $8\times 2=16$ 个实分量，
%  与 5d $\cN=1$（8 个超荷）一致；切勿与例外李代数 $\mathfrak{f}_4$（$\dim=52$、秩 4）混淆，
%  后者与超共形代数无关。此处照抄原文记号，仅提示读者。]

在本节中我们依次描述：对称代数 $\mathfrak{osp}(8^*|4)$（小节 2.1）、自对偶二形式规范场的性质（小节 2.2）、弦论／M 理论构造（小节 2.3），以及 $\cX(\fg)$ 的延展算子（extended operators，小节 2.4）。

\subsection{超共形代数}

维数 $d>2$ 的超共形代数已在某些附加条件下被 Nahm [154] 分类。它们的偶（玻色）部分由共形代数 $\fso(2,d)$（洛伦兹号差下）和一个 R 对称代数组成；奇（费米）部分由超荷组成，后者必须在 $\fso(2,d)$ 的旋量表示中变换，并且要使得平移由超荷的反交换子实现。

分类见表 2。在 $d=3,4,6$ 维，借助偶然同构\footnote{对 $d>6$ 不存在这类偶然同构，这或多或少解释了高维超共形代数的缺失。} $\fso(2,3)=\fsp(4,\R)$、$\fso(2,4)=\fsu(2,2)$ 与 $\fso(2,6)=\fso^*(8)$，共形代数与预期的 $\fso(2,d)$ 一致。在每种情形，$\fso(2,d)$ 的旋量表示恰好是另一个群的基本（矢量）表示。已知对 $d\ge 4$ 不存在庞加莱超荷多于 16 个的 SCFT（而 $d=3$ 时这类理论是自由的）[155]，这就给出表中对 $\cN$ 的上界。

% [Tutor Note: p.13「多于 16 个庞加莱超荷」这一上界的可核对之处：4d 的最大情形 $\cN=4$ 恰有 $4\times 4=16$
%  个超荷、6d $(2,0)$ 恰有 $2\times 8=16$ 个（6d 手征旋量 8 个实分量），故 "16" 正是 Nahm 表中
%  $\cN$ 上界的来源；3d 的 $\cN=8$ 有 $8\times 2=16$ 个（3d 旋量 2 个实分量）。一旦某项超共形代数的
%  超荷数 $>16$，其单粒子表示中必出现高自旋（$s>2$）态，而 [155] 排除了这类理论（3d 例外为自由场）。]

对于我们所关心的 6d 情形，洛伦兹代数 $\fso(2,6)$ 的最小旋量表示是手征的，而超共形代数包含 $\cN=1$ 或 $2$ 个手征性相同的手征旋量（技术上讲，是辛 Majorana--Weyl 旋量）。因此这些代数被称为 6d $\cN=(1,0)$ 与 6d $\cN=(2,0)$ 超共形代数。不存在 6d $\cN=(1,1)$ 超共形代数。我们关心的是其中最大的一个：6d $(2,0)$ 代数 $\mathfrak{osp}(8^*|4)$。

该代数的超荷在共形代数与 R 对称代数 $\fso(6,2)\times\fso(5)_R$ 的 $(\mathbf{8}_s,\mathbf{4})$ 表示\footnote{我们用粗体的维数来记单李代数的不可约表示（irrep）。出现歧义时有标准的修饰记号来区分，例如用上划线表示取共轭表示，用撇号表示存在若干同维数且彼此不共轭的不可约表示。一个特殊的例子是 $\fso(8)$ 及其它实形式如 $\fso(p,8-p)$：它们有三个 8 维不可约表示，即 $\fso(8)$ 的定义表示 $\mathbf{8}_v$，以及两个共轭的旋量表示 $\mathbf{8}_s$ 与 $\mathbf{8}_c$，后者由 $\fso(8)$ 的三重性（triality）自同构相联系。} 中变换，并带一个实条件。

% [原文 p.14]
把它分解到洛伦兹代数 $\fso(1,5)$ 的表示，得到 $(\mathbf{4},\mathbf{4})\oplus(\overline{\mathbf{4}},\mathbf{4})$，并带一个辛实条件。其中一组 $(\mathbf{4},\mathbf{4})$ 由庞加莱超荷组成，另一组是超共形变换。

\subsection{自对偶形式}

6d $\cN=(2,0)$ SCFT $\cX(\fg)$ 粗略地说是自对偶二形式规范场的理论，规范李代数 $\fg$ 取自 $\mathfrak{a}_{N-1}$、$\mathfrak{d}_N$、$\mathfrak{e}_6$、$\mathfrak{e}_7$、$\mathfrak{e}_8$，下面我们来解释。

\para{阿贝尔自对偶形式。}在偶数维 $d$ 中存在一个有意思的概念：（反）\footnote{自对偶与反自对偶情形只差一个符号，为简单起见我们一律写“自对偶”。}自对偶 $k$-形式，其中 $k=d/2-1$：即一个 $k$-形式 $B$，其分量 $B_{\alpha_1\dots\alpha_k}$（对 $\alpha_1,\dots,\alpha_k$ 反对称）满足：场强 $H=\dd B$ 在 Hodge 星下映为自身的倍数，即
\begin{equation}
H_{\alpha_0\alpha_1\dots\alpha_k} := (k+1)!\,\partial_{[\alpha_0}B_{\alpha_1\dots\alpha_k]} = \pm i^{d/2+s}\epsilon_{\alpha_0\dots\alpha_k\beta_0\dots\beta_k}\partial^{[\beta_0}B^{\beta_1\dots\beta_k]} .
\tag{2.1}
\end{equation}
这里方括号内的指标做反对称化，而 $i=\sqrt{-1}$ 的幂次中 $s=0$ 对应欧氏号差、$s=1$ 对应洛伦兹号差。自对偶条件针对的是场强，因此在规范变换 $B\to B+\dd\Lambda$（$\Lambda$ 为任意 $k$-形式）下不变：显式地说，这给 $k$-形式规范场 $B$ 的分量 $B_{\alpha_1\dots\alpha_k}$ 加上 $k!\,\partial_{[\alpha_1}\Lambda_{\alpha_2\dots\alpha_k]}$。

% [Tutor Note: p.14 式 (2.1) 之后那句现写作「加上 $k!\,\partial_{[\alpha_1}\Lambda_{\alpha_2\dots\alpha_k]}$」，
%  按 (2.1) 自身的方括号约定（即 $(k+1)!\,\partial_{[\ }B_{]}$ 给出无权重反对称和）反推，
%  正确的分量系数应是 $k$ 而不是 $k!$：因为 $\dd\Lambda$（$\Lambda$ 为 $(k-1)$-形式）的分量是
%  $(\dd\Lambda)_{\alpha_1\dots\alpha_k}=k\,\partial_{[\alpha_1}\Lambda_{\alpha_2\dots\alpha_k]}$，
%  由 $\delta H=\dd^2\Lambda=0$ 可知 $\delta B=\dd\Lambda$ 必须如此。逐例核对：$k=1$ 时 $1!=1$、
%  $k=2$ 时 $2!=2$，恰好与 $k$ 相同，故只有 $k\ge 3$ 才暴露差异（此时 $k!\ne k$）。
%  原文用 $k!$ 不影响"自对偶条件规范不变"这一结论（$\delta H$ 仍为零），只是分量系数写错。]

% [Tutor Note: p.14「实自对偶 $k$-形式只在 $d/2+s$ 为偶数时存在」的判据可直接由 (2.1) 读出：
%  实 $B$ ⇒ 实 $H$ ⇒ 右端的 $i^{d/2+s}$ 必须是实数，而 $i^n$ 为实当且仅当 $n$ 为偶数。
%  核对三个例子：4d 欧氏（$d/2+s=2$）给出 $i^2=-1$，即 $F=\mp\star F$ 的（反）自对偶瞬子；
%  6d 洛伦兹（$3+1=4$）给出 $i^4=1$，即实自对偶二形式；2d 洛伦兹（$1+1=2$）给出实标量
%  （只沿一个类光方向传播），与原文一致。]

由 (2.1) 可见，实自对偶 $k$-形式只在 $d/2+s$ 为偶数时存在。在 2d 这发生在洛伦兹号差下，它对应一个只沿一个类光方向传播的实标量场。（在欧氏情形，它是一个依赖单个全纯坐标的复手征玻色子。）在 4d 欧氏号差下，(2.1) 定义的是自对偶规范场位形，亦称瞬子，它们在 AGT 对应的 4d 一侧起关键作用。（在洛伦兹情形它们是复鞍点。）在 6d 洛伦兹号差下，我们得到实的自对偶二形式规范场 $B_{\alpha\beta}$。

我们关心 6d $(2,0)$ 超对称，此时含 $B_{\alpha\beta}$ 的多重态由 $B$、旋量 $\lambda$ 与标量 $\Phi$ 组成，它们分别按 R 对称群 $\mathfrak{usp}(4)=\fso(5)$ 的单态、4 维表示与 5 维表示变换。

\para{紧化在圆上。}让我们把这个 6d $(2,0)$ 阿贝尔 $(B,\lambda,\Phi)$ 理论放到圆上，并分解为 Kaluza--Klein（KK）模。如下面的习题所确定的那样，五个标量 $\Phi_I$ 仍是标量，旋量 $\lambda$ 也是，而自对偶二形式规范场 $B$ 变成 5d 的通常规范场 $A$。合起来就给出阿贝尔 5d $\cN=2$ SYM。

我们将在下面习题 2.1 中复习维数约化。从 $\cX(\fg)$ 约化到 5d 的一个重要方面是：5d SYM 有瞬子粒子，即在每个空间切片上具有非平凡拓扑数 $\int\epsilon^{\mu\nu\rho\sigma}F_{\mu\nu}F_{\rho\sigma}\dd^4x$ 的规范场位形。

% [Tutor Note: p.14 末的论断「瞬子粒子充当 KK 模塔」的量纲核对：5d 规范耦合的 $[g_{5d}^2]=$ 长度，
%  瞬子作用量 $8\pi^2/g_{5d}^2$ 的量纲是 $1/$长度，故瞬子质量 $\propto 1/g_{5d}^2$，正与半径 $R=g_{5d}^2$
%  的 KK 塔质量 $1/R$ 相符（原文 p.12 要点框已给出 $R=g_{5d}^2$）。这也是为什么 6d $(2,0)$ 理论
%  在圆上的 KK 模可以等价地用 5d 的瞬子（而非微扰 KK 激发）来组织。]

% [原文 p.15]
这些规范场 $A$ 的激发扮演 KK 模塔的角色：它们的质量（正比于）$1/g_{5d}^2$ 恰好被认定为 KK 模的质量 $1/R$。

\begin{exercise}{2.1}
\begin{enumerate}[label=\arabic*.,leftmargin=1.6em,itemsep=2pt,topsep=4pt]
\item 考虑一个 $D$ 维标量场 $\varphi$，拉格朗日量 $\cL(\varphi)=\partial_\alpha\varphi\partial^\alpha\varphi-V(\varphi)$（为简单起见可取 $V=0$）。把它放在 $d$ 维 Minkowski 空间乘以半径为 $R$ 的 $(D-d)$ 维环面上（为简单起见可取 $D-d=1$）。写出 $\varphi$ 沿圆方向的 Fourier 分解，并把 $\varphi$ 的作用量改写为这些分量的作用量。在 $R\to 0$ 的极限下，注意除零模外所有 Fourier 模都变成无穷大质量。
\item 对阿贝尔矢量场 $A_\alpha$（$\alpha=0,\dots,D-1$）重复上述习题，拉格朗日量为 $F_{\alpha\beta}F^{\alpha\beta}$，其中 $F_{\alpha\beta}=\partial_\alpha A_\beta-\partial_\beta A_\alpha$。验证维数约化后的理论既有矢量场 $A_\mu$（$\mu=0,\dots,d-1$）又有 $D-d$ 个标量场。对有限的 $R$，后者可以规范不变地理解为 $A_\alpha$ 沿环面各坐标圆的 Wilson 圈。$D$ 维的规范变换如何作用？由此推出这些标量场取值在圆上。
\item 对从 6d 约化到 5d 的二形式 $B_{\alpha\beta}$ 重复上述习题。结果得到一个二形式 $B_{\mu\nu}$ 和一个一形式 $A_\mu$。对 $B_{\alpha\beta}$ 施加自对偶条件，验证 $B_{\mu\nu}$ 可以由 $A_\mu$ 重构出来（至多差一个规范变换）。
\end{enumerate}
\end{exercise}

\para{非阿贝尔理论。}回顾 4d 中的 Bianchi 恒等式 $\partial_{[\mu}F_{\nu\rho]}=0$。它推广为 $\dd H=\dd\dd B=0$。对自对偶形式，这蕴含自由运动方程 $\dd\star\dd B=0$，即 $\partial^\mu H_{\mu\nu\dots}=0$。我们如何加入相互作用？在 4d，定义瞬子的方程 $F_{\mu\nu}=\mp\frac{1}{2}\epsilon_{\mu\nu\rho\sigma}F^{\rho\sigma}$ 对非阿贝尔规范场的场强也说得通：$F=\dd A+A\wedge A$，显式地 $F_{\mu\nu}=\partial_\mu A_\nu-\partial_\nu A_\mu+[A_\mu,A_\nu]$。Bianchi 恒等式的非阿贝尔版本是 $\epsilon^{\lambda\mu\nu\rho}D_\mu F_{\nu\rho}=0$。当规范场位形自对偶时，这蕴含标准的 Yang--Mills 运动方程 $D_\mu F^{\mu\lambda}=0$。相反，对其它自对偶 $k$-形式（$k\ne 1$），关系 $H=\dd B$ 没有明显的非阿贝尔推广，因而没有明显的引入相互作用的办法。于是我们改用两种弦论构造。

% [Tutor Note: p.15「$k\ne 1$ 没有明显的非阿贝尔推广」的实际障碍：$k=1$ 时 $F=\dd A+A\wedge A$ 中的
%  $A\wedge A$ 之所以可能，是因为 $A$ 是联络、其规范变换是 $A\to gAg^{-1}+g\,\dd g^{-1}$（仿射作用，
%  平移部分正是恰当形式）；$k\ge 2$ 的自对偶形式若要非阿贝尔化，需要更高层的 gerbe／2-群结构，
%  仅用一个 Lie 代数取值的 $k$-形式写不出自洽的 $H=\dd B+[B,B]$：核对方法是要求该推广的规范变换
%  下 $H$ 协变，而矩阵值 $k$-形式的楔积 $[B,B]$ 只在 $k$ 为奇数时才非零（$k=1$ 恰好如此）。]

\subsection{6d 理论的膜构造}

6d $(2,0)$ 理论 $\cX(\fg)$ 自然地作为某个 6d $(2,0)$ 小弦理论（little string theory）的零弦张力极限出现，后者的激发是自对偶弦。这些自对偶弦先以超引力解的形式被发现 [156]，随后出现在 IIB 弦论中的 ADE 轨形奇点处 [157]，再后来表现为终止于 NS5 膜上的 D2 膜（或终止于 M5 膜上的 M2 膜）[158]。（也存在类似的 6d $(1,1)$ 小弦理论，但它们没有共形极限，故我们不再讨论。）6d $(2,0)$ 小弦理论像 $\cX(\fg)$ 一样由单链\footnote{这里的 ADE 分类来自弦世界面上的反常抵消 [159]。}李代数 $\fg$ 标记，我们讨论 [160] 中综述的两种弦论构造。

\begin{itemize}[leftmargin=1.6em,itemsep=2pt,topsep=2pt]
\item 在 IIA 弦论中一叠 $N$ 个重合 NS5 膜（或在 M 理论中一叠 M5 膜）的零弦耦合极限 $g_s\to 0$ 下，体自由度退耦，得到 $\fg=\fsu(N)$ 的 6d $(2,0)$ 小弦理论。
\end{itemize}
